\section{In-Context Learning as Online Belief Propagation}
\label{sec:icl}

In-context learning (ICL) is the phenomenon that a transformer, given a
few examples in the prompt, performs a new task without any weight updates.
This is empirically reliable and practically important, but lacks a
consensus mechanistic explanation.

\subsection*{The BP Interpretation}

Under the BP interpretation, in-context learning is online Bayesian
updating. Each in-context example is additional evidence injected into the
implicit factor graph. The forward pass at each new token is one round of
message passing that conditions the current belief state on the accumulating
evidence. The context window is not just a memory buffer --- it is the
evidence base over which inference is performed.

This gives a natural account of several ICL phenomena. Performance improves
with more examples because more evidence tightens the posterior. Performance
improves with more layers because more layers means more rounds of BP,
allowing the evidence to propagate further through the factor graph. Tasks
that are easy for ICL are tasks where the relevant factor graph is simple
enough that a few examples suffice to identify the correct posterior.

Xie et al.~\cite{xie2022icl} propose a related Bayesian account of ICL at
the population level: the model maintains a posterior over latent document
concepts and updates it as it sees more examples. The BP account is the
mechanistic complement: each forward pass is one round of BP on the implicit
factor graph, and ICL works because the factor graph encodes the conditional
relationships learned from training.

\subsection*{What This Predicts}

If ICL is BP-based inference, then:

\begin{itemize}
\item ICL performance on a task should improve with model depth, up to the
point where the factor graph converges --- and then plateau.
\item On synthetic tasks with known factor graph structure, ICL performance
should track the convergence of BP on that graph.
\item Tasks requiring many inference hops should benefit more from depth
than tasks requiring few.
\item ICL should fail on tasks whose factor graph structure is not
representable within the model's implicit graph --- not just tasks that
are ``too hard.''
\end{itemize}

\subsection*{Epistemic Status}

\begin{center}
\begin{tabular}{ll}
\toprule
Claim & Status \\
\midrule
ICL is consistent with online Bayesian updating & Mechanistic interpretation \\
Each in-context example injects evidence into the factor graph & Mechanistic interpretation \\
ICL performance tracks BP convergence on synthetic graphs & Empirical hypothesis \\
Depth improves ICL via more inference rounds & Empirical hypothesis \\
Population-level Bayesian ICL account~\cite{xie2022icl} & External empirical result \\
\bottomrule
\end{tabular}
\end{center}